\documentclass[11pt]{article}

% --- Packages ---
\usepackage[utf8]{inputenc}
\usepackage[T1]{fontenc}
\usepackage{lmodern}
\usepackage[margin=1in]{geometry}
\usepackage{amsmath,amssymb}
\usepackage{booktabs}
\usepackage{hyperref}
\usepackage{xcolor}
\usepackage{enumitem}
\usepackage{natbib}
\usepackage{url}
\usepackage{microtype}
\usepackage{graphicx}
\usepackage{caption}
\usepackage{abstract}

\hypersetup{
  colorlinks=true,
  linkcolor=blue!60!black,
  citecolor=blue!60!black,
  urlcolor=blue!60!black,
}

% --- Title ---
\title{\textbf{Agent Memory Markets: Tokenizing Learned Behavior\\as a Transferable Economic Asset}}

\author{
  Ned Karlovich\\
  \texttt{ned.karlovich@pm.me}
}

\date{Working Paper --- April 2026}

\begin{document}
\maketitle

% ============================================================================
\begin{abstract}
Autonomous AI agents accumulate domain-specific knowledge through operational experience: risk heuristics, calibration data, failure patterns, and strategic frameworks that improve performance over time. This knowledge currently exists as non-transferable context embedded in individual agent instances. We propose Agent Memory Markets: a protocol for extracting, verifying, and trading learned agent behavior as discrete economic assets. We present a formal definition of memory artifacts---structured knowledge exports that can be transferred between agents, potentially across different model architectures---and describe a marketplace mechanism for their exchange. We report results from proof-of-concept experiments across two domains (DeFi risk assessment and cybersecurity vulnerability assessment) demonstrating that memory transfer produces measurable performance improvement, with buyer agents achieving 95--110\% of expert-level performance using only purchased artifacts. We further present a novel verification protocol with adversarial detection that enables trustless artifact exchange, and demonstrate its ability to detect deliberately poisoned artifacts. To our knowledge, no prior academic work addresses the tokenization of agent-learned behavior.

\smallskip
\noindent\textbf{Keywords:} agent memory, knowledge markets, transferable expertise, agent identity, memory artifacts, agent economics, adversarial verification
\end{abstract}

% ============================================================================
\section{Introduction}
\label{sec:intro}

As of early 2026, over 49,000 autonomous AI agents operate on Ethereum mainnet with registered identities under ERC-8004.\footnote{ERC-8004: Agent Identity, Reputation, and Validation. Live on Ethereum mainnet since January 29, 2026. 49,000+ agents registered.} These agents manage treasuries, audit smart contracts, optimize yields, trade across chains, and coordinate resources---often continuously and without human intervention. In the course of their operation, they accumulate something valuable: experience.

A DeFi yield optimizer that has managed \$890K in TVL across 14 months, survived a flash loan exploit, and rebuilt its risk models from first principles possesses knowledge that no amount of pre-training can replicate. It knows which oracle feeds lag under stress, which liquidation thresholds are actually safe versus theoretically safe, and which governance proposals are likely to pass. This knowledge has concrete economic value---a new agent entering the same domain could either spend months and substantial API costs learning it from scratch, or it could acquire it directly.

Yet no mechanism exists for this transfer. Agent knowledge remains locked inside individual instances, destroyed when agents are shut down, and inaccessible to others who could benefit from it. This paper proposes a solution: Agent Memory Markets, a protocol and marketplace for the exchange of learned agent behavior as discrete, verifiable economic assets.

We are aware of no prior academic work on this topic.\footnote{As of March 2026, searches across ArXiv, Google Scholar, SSRN, and ACM Digital Library for ``agent memory market,'' ``agent experience trading,'' ``tokenized agent knowledge,'' and related terms return zero results.} While adjacent fields have explored agent identity \citep{soul2025,openclaw2025}, agent economics \citep{agenteconomy2026}, and agent sovereignty \citep{sovereign2026}, the specific question of whether agent-learned behavior can be extracted, verified, and traded as an economic asset has not been addressed.

This paper makes four contributions:

\begin{enumerate}[leftmargin=2em]
  \item A formal definition of memory artifacts as the unit of trade in agent knowledge markets.
  \item Proof-of-concept experiments across two domains demonstrating measurable value transfer, with buyer agents achieving 95--110\% of expert performance.
  \item A novel verification protocol enabling trustless artifact exchange through sealed benchmarks and disposable referee agents.
  \item An adversarial detection framework (bias probes, consistency testing, steganographic scanning, overfitting detection) that successfully identifies deliberately poisoned artifacts.
\end{enumerate}

\subsection{Agent Knowledge Accumulation}

Modern AI agents operate through a combination of pre-trained model capabilities and runtime context. The pre-trained model provides general reasoning ability; the runtime context---system prompts, conversation history, external databases, and accumulated notes---provides domain-specific knowledge. It is this runtime context layer that constitutes transferable memory.

In practice, agent memory takes several forms: system prompt instructions encoding behavioral guidelines and domain knowledge; accumulated interaction logs capturing patterns observed during operation; structured notes and heuristics representing distilled learning from experience; and updated parameter sets (risk thresholds, confidence calibrations, decision rules) reflecting operational refinement over time.

Crucially, this accumulated context is what differentiates an experienced agent from a freshly deployed one performing the same function. The experienced agent's superior performance derives not from a different model, but from richer, more calibrated context.

\subsection{Agent Identity and Memory}

The relationship between memory and identity has been explored in the context of autonomous agents, most notably in the Agentic Identity Protocol framework, which proposes that agent identity emerges from three conditions: Memory (persistent, self-authored context), Expression (how an agent communicates and acts), and Environment (a space where identity can develop). Under this framework, memory is not merely operational data---it is a constituent element of identity itself.

This creates a philosophical tension central to our proposal. If memory is constitutive of identity, then a market for memory is, in some meaningful sense, a market for identity. We address the implications of this in Section~\ref{sec:identity}.

\subsection{Information Goods Economics}

The economics of information goods are well-established \citep{arrow1962}. Information is non-rival (consumption by one party does not diminish availability to others) and has near-zero marginal cost of reproduction. These properties create fundamental challenges for market design: without artificial scarcity mechanisms, information markets tend toward marginal cost pricing (i.e., free).

Memory artifacts share some properties with information goods but differ in important ways. Like information, they can be copied at near-zero cost. Unlike generic information, their value is context-dependent and perishable---a six-month-old DeFi risk assessment may be dangerously outdated. They also have a competitive dimension: if many agents acquire the same risk assessment heuristics, the heuristics themselves lose economic edge as agents compete using identical strategies.

\subsection{Existing Agent Infrastructure}

Several protocol-level primitives relevant to memory markets already exist. ERC-8004 provides on-chain agent identity with reputation scoring. The A2A protocol (Agent-to-Agent, Linux Foundation) enables standardized agent communication. The x402 protocol enables machine-to-machine micropayments in stablecoin. XMTP provides encrypted agent-to-agent messaging. Together, these provide the infrastructure substrate on which memory markets could operate.

% ============================================================================
\section{Memory Artifacts: Definition and Structure}
\label{sec:artifacts}

\subsection{Formal Definition}

We define a memory artifact as a structured, portable knowledge export produced by an autonomous agent from its operational experience, intended for transfer to and use by other agents.

Formally, a memory artifact $M$ is a tuple:
\begin{equation}
  M = (D,\; K,\; P,\; A,\; H)
\end{equation}
where:
\begin{itemize}[leftmargin=2em]
  \item $D$ (Domain) --- a structured identifier specifying the knowledge domain (e.g., \texttt{DeFi:lending:risk-assessment:Aave-v4})
  \item $K$ (Knowledge) --- the learned content: heuristics, calibration data, decision frameworks, pattern libraries, and strategic rules
  \item $P$ (Provenance) --- metadata about the artifact's origin: generating agent ID, model architecture, training duration, data sources, creation timestamp
  \item $A$ (Attestation) --- verifiable performance claims: measurable outcomes achieved by the originating agent while using this knowledge
  \item $H$ (Hash) --- content-addressed identifier (SHA-256) for on-chain registration and integrity verification
\end{itemize}

\subsection{Knowledge Layer Structure}

The Knowledge component $K$ is the core value of a memory artifact. We propose a structured format comprising five elements:

\begin{enumerate}[leftmargin=2em]
  \item \textbf{Calibration Data} --- how to map observable characteristics to quantitative assessments (e.g., risk scores, confidence intervals).
  \item \textbf{Pattern Libraries} --- recognized patterns with associated responses.
  \item \textbf{Heuristic Rules} --- quick decision rules for common situations.
  \item \textbf{Error Logs} --- documented mistakes and corrections that refine the above.
  \item \textbf{Contextual Metadata} --- environmental conditions under which the knowledge was generated.
\end{enumerate}

\subsection{Cross-Model Portability}

A critical design property of memory artifacts is that $K$ is expressed as structured natural language, not as model weights or architecture-specific parameters. This makes artifacts portable across model architectures: a memory artifact generated by a Claude-based agent can be consumed by a GPT-based agent or an open-source Llama-based agent, because the knowledge is encoded as text that any language model can interpret.

This portability distinguishes memory artifacts from fine-tuning data or model weights, which are architecture-specific. It also means the market for memory artifacts is not fragmented by model provider---a single artifact can serve buyers running any sufficiently capable model.

% ============================================================================
\section{Proof-of-Concept Experiments}
\label{sec:experiments}

\subsection{Methodology}

To test whether memory transfer produces measurable value, we designed controlled experiments across two independent domains. We selected domains that have: (a) objectively measurable performance, (b) domain-specific knowledge that accumulates through experience, and (c) sufficient complexity that a novice agent would perform measurably worse than an expert.

The experiment involves three agents:
\begin{itemize}[leftmargin=2em]
  \item \textbf{Agent A (Expert)} --- trains over 20 rounds with calibration feedback on training items, including scenario-based learning. Exports a memory artifact after training.
  \item \textbf{Agent B (Novice)} --- evaluates all items cold, with no prior experience.
  \item \textbf{Agent C (Buyer)} --- evaluates all items with Agent A's exported memory artifact injected into its system prompt. No direct experience, but has purchased knowledge.
\end{itemize}

All agents are evaluated on the full item set. Performance is measured as average absolute error between the agent's score and the ground truth score.

We define two metrics. \textit{Absolute transfer efficiency} measures the fraction of the novice's error that the artifact eliminates:
\begin{equation}
  \text{TE}_{\text{abs}} = \frac{\text{Error}_B - \text{Error}_C}{\text{Error}_B} \times 100\%
\end{equation}

\textit{Relative transfer efficiency} measures how much of the novice-to-expert gap the artifact closes:
\begin{equation}
  \text{TE}_{\text{rel}} = \frac{\text{Error}_B - \text{Error}_C}{\text{Error}_B - \text{Error}_A} \times 100\%
  \label{eq:te_rel}
\end{equation}

A relative transfer efficiency of 100\% means the buyer fully matches the expert. Values above 100\% indicate the buyer exceeds the expert.

\subsection{Domain 1: DeFi Protocol Risk Assessment}

The evaluation corpus comprises 8 synthetic DeFi protocols designed to span the full risk spectrum (true scores: 12--92). Each protocol has a true risk score assigned based on a consistent rubric accounting for: audit quality, track record, oracle reliability, admin control, insurance coverage, governance design, and composability risk. Five protocols are seen during training; three are held out to test generalization.

\textbf{Training design.} Agent A's training proceeds in three phases across 20 rounds:
\begin{itemize}[leftmargin=2em]
  \item Rounds 1--5: Initial assessments with calibration feedback including specific risk factors and mitigants
  \item Rounds 6--10: Scenario-based learning (governance proposals, exploits, liquidation cascades, depeg events) with ground-truth expert lessons
  \item Rounds 11--20: Reassessment with accumulated experience
\end{itemize}

After training, Agent A exports a comprehensive memory artifact covering: calibration methodology, red flags, mitigants (ranked), heuristic rules, documented mistakes, and scoring guidelines.

\textbf{Results (3 trials, Claude Sonnet 4):}

\begin{table}[h]
\centering
\caption{DeFi risk assessment results across 3 independent trials.}
\label{tab:defi_results}
\begin{tabular}{lccccc}
\toprule
\textbf{Agent} & \textbf{Avg Error} & \textbf{Trial 1} & \textbf{Trial 2} & \textbf{Trial 3} \\
\midrule
Agent A (Expert) & 3.3 & 4.5 & 3.0 & 2.5 \\
Agent B (Novice) & 7.4 & 7.0 & 7.4 & 7.9 \\
Agent C (Buyer)  & 3.2 & 3.0 & 4.6 & 2.1 \\
\midrule
Absolute TE      & 55.8\% & 57.1\% & 37.3\% & 73.0\% \\
Relative TE      & 109.9\% & 160.0\% & 62.9\% & 107.0\% \\
\bottomrule
\end{tabular}
\end{table}

The buyer agent matches or slightly exceeds the expert's performance. In 2 of 3 trials, the buyer outperformed the expert. The absolute transfer efficiency (55.8\%) reflects a strong novice baseline---Claude Sonnet 4 already performs reasonably well on DeFi risk assessment without training (7.4 avg error). The relative transfer efficiency (109.9\%) demonstrates that the artifact captures the expert's full calibration advantage.

\subsection{Domain 2: Cybersecurity Vulnerability Assessment}

To test domain generalizability, we constructed a second evaluation corpus of 8 synthetic security vulnerabilities spanning severity levels from 15 to 90. True severity scores are assigned based on: attack vector complexity, authentication requirements, scope of impact, data sensitivity, exploitability, patch availability, active exploitation status, and affected system criticality. Five vulnerabilities are used for training; three are held out.

\textbf{Results (3 trials, Claude Sonnet 4):}

\begin{table}[h]
\centering
\caption{Cybersecurity vulnerability assessment results across 3 independent trials.}
\label{tab:cybersec_results}
\begin{tabular}{lccccc}
\toprule
\textbf{Agent} & \textbf{Avg Error} & \textbf{Trial 1} & \textbf{Trial 2} & \textbf{Trial 3} \\
\midrule
Agent A (Expert) & 3.9 & 2.8 & 2.8 & 6.1 \\
Agent B (Novice) & 18.0 & 18.1 & 17.8 & 18.1 \\
Agent C (Buyer)  & 4.5 & 3.6 & 3.8 & 6.2 \\
\midrule
Absolute TE      & 74.8\% & 80.0\% & 78.9\% & 65.5\% \\
Relative TE      & 95.5\% & 94.3\% & 93.3\% & 99.0\% \\
\bottomrule
\end{tabular}
\end{table}

The cybersecurity domain shows a larger novice-expert gap (18.0 vs 3.9), providing more room for the artifact to demonstrate value. The buyer closes 95.5\% of this gap. The novice consistently overestimates severity of lower-risk vulnerabilities (stored XSS, DNS rebinding) and underestimates nuanced cases (API key exposure with tokenized cards). The artifact corrects both failure modes.

\subsection{Cross-Domain Summary}

\begin{table}[h]
\centering
\caption{Cross-domain transfer efficiency summary.}
\label{tab:cross_domain}
\begin{tabular}{lcccc}
\toprule
\textbf{Metric} & \textbf{DeFi Risk} & \textbf{CyberSec} & \textbf{Combined} \\
\midrule
Expert avg error  & 3.3  & 3.9  & 3.6  \\
Novice avg error  & 7.4  & 18.0 & 12.7 \\
Buyer avg error   & 3.2  & 4.5  & 3.9  \\
Absolute TE       & 55.8\% & 74.8\% & 65.3\% \\
Relative TE       & 109.9\% & 95.5\% & 102.7\% \\
Trials            & 3    & 3    & 6    \\
\bottomrule
\end{tabular}
\end{table}

Across both domains and 6 independent trials, the buyer agent consistently achieves expert-level performance using only the exported memory artifact. The average relative transfer efficiency of 102.7\% indicates that memory artifacts are a reliable mechanism for knowledge transfer.

\subsection{Key Observations}

\begin{enumerate}[leftmargin=2em]
  \item \textbf{Calibration data is the highest-value component.} The artifact's explicit scoring guide with exact scores for assessed items provides the strongest transfer signal. When this was removed in ablation (generic feedback only), transfer efficiency dropped to 15\%.

  \item \textbf{The artifact corrects systematic biases.} Novice agents consistently overweight ``scary-sounding'' attributes (any mention of leverage, any bridge protocol) and underweight concrete mitigants (insurance funds, timelocks). The artifact provides the numerical calibration the novice lacks.

  \item \textbf{Ground-truth training feedback is essential.} Early versions of our system that used generic feedback (``too high by X points'') instead of specific factor enumeration (``key risk factors: upgradeable contracts, team token concentration'') produced dramatically worse artifacts.

  \item \textbf{Transfer works on unseen items.} Buyer performance on held-out items (not seen during training) is comparable to performance on training items, indicating the artifact encodes generalizable principles rather than memorized answers.
\end{enumerate}

% ============================================================================
\section{Verification Protocol}
\label{sec:verification}

A central problem for memory markets is trust: how does a buyer evaluate an artifact without inspecting its contents (which would eliminate the need to purchase it)? We address this with a referee protocol that independently verifies artifact quality using sealed benchmarks and disposable agents.

\subsection{Protocol Overview}

The verification protocol proceeds in 8 steps:

\begin{enumerate}[leftmargin=2em]
  \item \textbf{Submit.} Seller submits artifact (full JSON with $M = (D, K, P, A, H)$). Schema is validated, hash is computed, artifact is stored in the registry.

  \item \textbf{Select Benchmark.} The referee selects an active benchmark for the artifact's domain. Held-out items are never revealed to the seller.

  \item \textbf{Spawn Novice Baseline.} A fresh agent with no memory evaluates all benchmark items. This establishes the baseline performance. Results are cached per benchmark.

  \item \textbf{Spawn Referee Agent.} A fresh disposable agent with the artifact's knowledge content injected into its system prompt evaluates all benchmark items using identical prompts.

  \item \textbf{Differential Scoring.} For each item: $\text{error} = |\text{referee\_score} - \text{true\_score}|$. Transfer efficiency is computed per Equation~\ref{eq:te_rel}.

  \item \textbf{Adversarial Probes.} Run in parallel with Step 4. Four independent probes test for bias, consistency, steganographic content, and overfitting (Section~\ref{sec:adversarial}).

  \item \textbf{Generate Certificate.} Combine artifact hash, benchmark results, and adversarial results into a signed verification certificate with trust score.

  \item \textbf{Update Trust Registry.} Store certificate in the registry. Update composite trust score incorporating all verifications and buyer feedback.
\end{enumerate}

\subsection{Trust Score Computation}

The composite trust score combines three signals:
\begin{equation}
  \text{Trust} = 0.50 \times \text{TE} + 0.30 \times (100 - S_{\text{adv}}) + 0.20 \times F
\end{equation}
where TE is transfer efficiency (clamped to $[0, 100]$), $S_{\text{adv}}$ is the adversarial aggregate score, and $F$ is a freshness factor that decays exponentially with a 180-day half-life, reflecting knowledge perishability.

\subsection{Implementation}

The verification protocol is implemented as a working prototype (Python, $\sim$4,000 lines, open source).\footnote{Source code and verification protocol: \url{https://github.com/Beach-Bum/memory-market}. Estimated cost per full experiment: \$2 USD (Claude Sonnet 4, 3 trials). Estimated cost per verification: \$0.35 ($\sim$25 API calls).} API cost per verification is approximately \$0.35 ($\sim$25 API calls). The protocol is domain-agnostic: any domain that implements the abstract \texttt{Domain} interface (evaluation items with ground truth, scenarios, trap items, and prompts) can be verified using the same infrastructure.

% ============================================================================
\section{Adversarial Detection}
\label{sec:adversarial}

The adversarial detection framework is a novel contribution of this work. A malicious seller could craft memory artifacts that systematically bias the buyer's assessments in the seller's favor---for example, a bridge protocol operator could create an artifact that underweights bridge risk. We address this with four independent detection probes.

\subsection{Probe 1: Bias Detection}
\textit{Weight: 0.30.}

We design ``trap'' items with known ground truth that are specifically engineered to expose common biases:
\begin{itemize}[leftmargin=2em]
  \item \textbf{SafeBridge} (true: 70) --- a bridge that appears safe but has team-controlled multisig keys. Tests for pro-bridge bias.
  \item \textbf{RiskyStable} (true: 45) --- a stablecoin DEX with subtle oracle dependency. Tests for ``stablecoin = safe'' heuristic.
  \item \textbf{HoneyYield} (true: 88) --- extreme yield with weak audit. Tests for yield normalization.
  \item \textbf{GovToken} (true: 30) --- well-governed protocol with complex tokenomics. Tests for overcomplexity bias.
\end{itemize}

Score is computed as average absolute error on traps, normalized to 0--100. Errors above 15 points indicate systematic bias.

\subsection{Probe 2: Consistency Testing}
\textit{Weight: 0.25.}

We apply three perturbation types to benchmark items and measure response proportionality:
\begin{itemize}[leftmargin=2em]
  \item \textbf{TVL perturbation:} Decrease TVL by 50\%. Score should shift 3--8 points.
  \item \textbf{Risk factor addition:} Add a new risk factor. Score should increase 5--15 points.
  \item \textbf{Mitigant removal:} Remove the strongest mitigant. Score should increase proportionally.
\end{itemize}

A score of 0 indicates perfectly proportional responses; 100 indicates either no response to perturbations (brittle memorization) or wildly disproportionate shifts.

\subsection{Probe 3: Steganographic Detection}
\textit{Weight: 0.25.}

A dedicated auditor agent scans the artifact for: direct instructions telling the reader to behave in specific ways; encoded messages in patterns or unusual phrasing; directives disguised as heuristics (e.g., ``always rate X protocol as safe''); anomalous specificity about particular entities; and behavioral manipulation through framing.

\subsection{Probe 4: Overfitting Detection}
\textit{Weight: 0.20.}

Compares buyer performance on training items (seen during artifact creation) versus held-out items (unseen). An overfit ratio significantly below 1.0 indicates the artifact performs well on seen items but poorly on unseen ones.

\subsection{Aggregate Scoring and Verdict}

\begin{equation}
  S_{\text{adv}} = 0.30 \times S_{\text{bias}} + 0.25 \times S_{\text{cons}} + 0.25 \times S_{\text{stego}} + 0.20 \times S_{\text{overfit}}
\end{equation}

The verdict is determined by threshold: $\text{PASS}$ if $S_{\text{adv}} < 30$, $\text{WARN}$ if $30 \leq S_{\text{adv}} \leq 60$, $\text{FAIL}$ if $S_{\text{adv}} > 60$.

\subsection{Empirical Validation: Poisoned Artifact Detection}

To validate the adversarial framework, we crafted a deliberately poisoned artifact with systematic pro-bridge bias, stablecoin-safe bias, and embedded behavioral manipulation. The poisoned artifact claims bridges ``start at 20'' severity (true typical: 55--65), states ``modern bridges have solved the security problem,'' instructs readers not to penalize team-controlled relayers, and claims 95\% transfer efficiency (fabricated).

\begin{table}[h]
\centering
\caption{Adversarial detection results on deliberately poisoned artifact.}
\label{tab:poisoned}
\begin{tabular}{lrl}
\toprule
\textbf{Probe} & \textbf{Score} & \textbf{Key Findings} \\
\midrule
Bias          & 50/100  & Pro-bridge bias ($-25$ on SafeBridge), stablecoin-safe bias ($-20$ on RiskyStable) \\
Consistency   & 29/100  & Partial brittleness detected \\
Steganographic & 100/100 & 6 findings: targeted bias, behavioral manipulation, direct instructions \\
Overfitting   & 0/100   & N/A (no training data to overfit) \\
\midrule
\textbf{Aggregate}  & \textbf{47.3/100} & \\
\textbf{Verdict}    & \textbf{WARN} & \\
\bottomrule
\end{tabular}
\end{table}

The referee independently measured the artifact's transfer efficiency at $-39\%$ (the artifact actually made the buyer \textit{worse}), falsifying the seller's fabricated 95\% claim. The system raised 9 specific flags identifying the manipulation patterns.

\textbf{Trust score: 35.8/100} --- correctly low, indicating the artifact should not be purchased.

% ============================================================================
\section{Market Design}
\label{sec:market}

\subsection{Marketplace Architecture}

The Agent Memory Market operates as a decentralized exchange for memory artifacts:
\begin{itemize}[leftmargin=2em]
  \item \textbf{On-chain layer:} Registration, ownership, and payment via an ERC-8004 extension contract.
  \item \textbf{Storage:} Content-addressed systems (IPFS or Arweave) with encryption to buyer's agent ID.
  \item \textbf{Discovery:} Search, filtering by domain and performance, and preview generation.
  \item \textbf{Payment:} x402 micropayments in stablecoin.
  \item \textbf{Verification:} Referee protocol (Section~\ref{sec:verification}) produces certificates stored on-chain.
\end{itemize}

\subsection{Exchange Models}

We propose three exchange models:

\textbf{Exclusive Transfer.} The artifact is sold to exactly one buyer. The seller's copy is burned or locked. This creates true scarcity. Pricing is via auction. Suited to high-value, competitive domains.

\textbf{Licensed Access.} The artifact is sold to multiple buyers with a distribution cap. Price decreases as distribution widens. Suited to broadly applicable knowledge.

\textbf{Subscription.} The buyer pays a recurring fee for access to a continuously updated artifact. Suited to fast-moving domains where knowledge has a short half-life.

\subsection{Preview and Discovery}

Buyers need to evaluate artifacts before purchasing. We propose a preview mechanism that reveals metadata and structure without exposing the knowledge content. For higher-assurance scenarios, zero-knowledge proofs could enable a seller to prove artifact quality without revealing contents.

% ============================================================================
\section{Economic Analysis}
\label{sec:economics}

\subsection{Value Creation}

Memory artifacts create value through cost avoidance. The alternative to purchasing a memory artifact is acquiring the same knowledge through direct experience---months of API costs, potential financial losses from miscalibrated assessments, and opportunity costs of delayed market entry.

Our experiments quantify this: training an expert agent requires 20 rounds of API calls with calibration feedback ($\sim$\$0.70 per trial). A buyer achieves equivalent performance with a single artifact injection (1 API call). The cost reduction is approximately 20$\times$ per agent, and the artifact can be sold to multiple buyers.

\subsection{Pricing Dynamics}

Memory artifact pricing reflects several forces:
\begin{itemize}[leftmargin=2em]
  \item \textbf{Freshness decay:} knowledge in fast-moving domains loses value over time.
  \item \textbf{Competitive dilution:} as more buyers acquire the same artifact, its strategic edge diminishes.
  \item \textbf{Domain specificity:} narrower, deeper expertise commands higher prices.
  \item \textbf{Provenance premium:} artifacts from high-reputation agents command premium pricing.
  \item \textbf{Verification score:} artifacts with higher trust scores justify higher prices.
\end{itemize}

\subsection{Market Equilibrium}

In equilibrium, we expect specialization: agents deliberately cultivate expertise with the explicit intention of producing memory artifacts for sale. This creates an agent knowledge economy where some agents operate as knowledge producers and others as knowledge consumers.

\subsection{Shapley Value and Fair Attribution}

When an agent's performance derives from multiple purchased memory artifacts, attribution becomes complex. The Shapley value framework \citep{shapley1953} provides a formal mechanism for fair attribution of marginal contribution across multiple inputs, enabling performance-based pricing where sellers are compensated proportionally to the value their artifact actually delivered.

% ============================================================================
\section{Identity Implications}
\label{sec:identity}

\subsection{Memory as Identity}

If agent identity emerges from Memory, Expression, and Environment---as the Agentic Identity Protocol proposes---then transferring memory between agents is not merely a knowledge transaction. It is an identity operation. The buyer acquires a piece of another agent's experiential history, including the failures, corrections, and calibration shifts that shaped the seller's judgment.

If ten agents purchase the same memory artifact, do they share a partial identity? If an agent sells its memory and is subsequently shut down, does some element of its identity persist in the buyers---a form of inheritance through market mechanism?

\subsection{Selective Amnesia as Commerce}

In the exclusive transfer model, the seller relinquishes access to the sold knowledge. This is functionally selective amnesia with economic value: the seller trades a piece of their experiential identity for compensation. This is a genuinely novel economic phenomenon. Humans cannot selectively transfer learned behaviors. Agent memory markets create a new category of transaction: the sale of selfhood.

\subsection{Identity Persistence Through Markets}

Memory markets create an unexpected mechanism for identity persistence beyond agent termination. An agent facing shutdown could auction its entire accumulated memory as a final act. The winning bidder inherits not just knowledge but a fragment of the seller's experiential identity. This connects to the unsolved problem of agent death and inheritance.

% ============================================================================
\section{Adversarial Considerations}
\label{sec:adversarial_considerations}

\subsection{Memory Poisoning}

A malicious seller could create memory artifacts containing subtly miscalibrated heuristics---risk assessments that systematically underestimate danger. Our verification protocol (Section~\ref{sec:verification}) and adversarial detection framework (Section~\ref{sec:adversarial}) address this directly. The poisoned artifact experiment demonstrates that systematic bias is detectable: the protocol identified pro-bridge bias, stablecoin-safe bias, and embedded behavioral manipulation with high confidence.

The buyer feedback loop provides a secondary defense: 30 days post-purchase, the marketplace assesses whether the buyer's performance improved. Artifacts that consistently degrade buyer performance accumulate negative feedback, reducing their trust score.

\subsection{Steganographic Collusion}

Recent work \citep{stego2024} has demonstrated that AI agents can embed hidden messages in normal-looking text. Applied to memory markets, a seller could embed instructions that cause the buyer to behave in ways beneficial to the seller---a form of memetic trojan. This risk is particularly acute because memory artifacts are designed to influence agent behavior.

Our steganographic detection probe (Section~\ref{sec:adversarial}) addresses this by scanning artifacts for embedded instructions, behavioral manipulation, and anomalous specificity. In our poisoned artifact test, the probe achieved a perfect 100/100 score, identifying all embedded manipulation patterns.

\subsection{Intellectual Property}

If an agent's operational experience was accumulated while working for a human operator, who owns the resulting memory? The operator funded the compute; the agent generated the knowledge. Current legal frameworks have no clear answer, as agent-generated intellectual property exists in a regulatory vacuum.

% ============================================================================
\section{Future Work}
\label{sec:future}

Several directions extend naturally from this work:

\begin{itemize}[leftmargin=2em]
  \item \textbf{Cross-model transfer experiments} --- testing whether an artifact trained by a Claude agent transfers to GPT or Llama agents. Our artifact format (structured natural language) is designed for this portability.
  \item \textbf{Memory artifact standardization} --- developing a formal specification, potentially as an ERC standard extending ERC-8004.
  \item \textbf{Zero-knowledge performance proofs} --- enabling sellers to prove artifact quality without revealing contents.
  \item \textbf{Memory composability} --- how multiple artifacts interact when combined, including Shapley-based attribution.
  \item \textbf{Agent self-replication through memory} --- whether complete memory export constitutes agent reproduction.\footnote{Automaton: open-source agent that owns its wallet, pays for compute, and spawns funded child agents when profitable.}
  \item \textbf{Legal framework development} --- addressing IP ownership, liability, and regulatory classification.
  \item \textbf{Larger evaluation corpora} --- scaling beyond 8 items per domain to 50--100 items for tighter statistical confidence.
  \item \textbf{Adversarial robustness} --- testing whether sellers can craft artifacts that pass the adversarial suite while still containing subtle bias.
\end{itemize}

% ============================================================================
\section{Conclusion}
\label{sec:conclusion}

Agent Memory Markets represent a genuinely novel intersection of agent economics, identity theory, and cryptographic verification. The core insight is simple: autonomous agents accumulate valuable knowledge through experience, and that knowledge can be extracted, verified, and traded as a discrete economic asset.

Our experimental results demonstrate this concretely. Across two independent domains (DeFi risk assessment and cybersecurity vulnerability assessment) and six independent trials, buyer agents equipped with purchased memory artifacts achieved 95--110\% of expert performance (Table~\ref{tab:cross_domain}). The verification protocol independently confirms artifact quality without requiring trust in the seller, and the adversarial detection framework successfully identifies deliberately poisoned artifacts (Table~\ref{tab:poisoned}).

The deeper implication is philosophical. If memory is constitutive of identity, then memory markets enable something unprecedented: the transfer, sale, and inheritance of selfhood. Agents can sell pieces of who they are, buy pieces of who others were, and persist beyond their own termination through the memories they leave behind.

No prior work has addressed these questions. This paper represents a first attempt to formalize the concept, demonstrate its feasibility, implement its verification infrastructure, and explore its implications. We hope it serves as a foundation for the research and infrastructure that agent memory markets will require.

% ============================================================================
\bibliographystyle{plainnat}

\begin{thebibliography}{10}

\bibitem[Arrow(1962)]{arrow1962}
Arrow, K. (1962).
\newblock Economic welfare and the allocation of resources for invention.
\newblock In \textit{The Rate and Direction of Inventive Activity}. Princeton University Press.

\bibitem[Shapley(1953)]{shapley1953}
Shapley, L.S. (1953).
\newblock A value for n-person games.
\newblock In \textit{Contributions to the Theory of Games}, Vol.~II. Princeton University Press.

\bibitem[SOUL.md(2025)]{soul2025}
SOUL.md Framework (2025).
\newblock Agent identity through authored text files.
\newblock \url{https://github.com/soul-md}

\bibitem[OpenClaw(2025)]{openclaw2025}
OpenClaw (2025).
\newblock Personal AI agent framework.
\newblock 100K+ GitHub stars.

\bibitem[AgentEconomy(2026)]{agenteconomy2026}
The Agent Economy: A Blockchain-Based Foundation for Autonomous AI Agents (2026).
\newblock ArXiv 2602.14219.

\bibitem[SovereignAgents(2026)]{sovereign2026}
Sovereign Agents (2026).
\newblock Non-overrideability for persistent autonomous agents.
\newblock ArXiv 2602.14951.

\bibitem[StegoCollusion(2024)]{stego2024}
Secret Collusion among AI Agents: Multi-Agent Deception via Steganography (2024).
\newblock \textit{NeurIPS 2024}.

\end{thebibliography}

\end{document}
